% 系外卫星（综述）
% license CCBYSA3
% type Wiki

本文根据 CC-BY-SA 协议转载翻译自维基百科\href{https://en.wikipedia.org/wiki/Exomoon}{相关文章}。

\begin{figure}[ht]
\centering
\includegraphics[width=8cm]{./figures/f73ad0be1621979a.png}
\caption{候选系外卫星 Kepler-1625b I 环绕其母行星运行的艺术想象图 $^{[1]}$} \label{fig_Exomoo_1}
\end{figure}
系外卫星（exomoon），或称太阳系外卫星（extrasolar moon），是指环绕系外行星或其他非恒星的太阳系外天体运行的天然卫星。$^{[2]}$

受限于现有观测技术，系外卫星的探测与确认十分困难，$^{[3]}$ 截至目前尚无被正式确认的系外卫星发现。$^{[4]}$ 不过，诸如开普勒（Kepler）任务的观测已经发现了若干候选体。$^{[5][6]}$ 此外，通过微引力透镜方法，还探测到两个可能环绕流浪行星运行的系外卫星候选体。$^{[7][8]}$ 2019 年 9 月，天文学家报告称，对“塔比星”（Tabby’s Star）观测到的变暗现象，可能是由一颗孤立系外卫星被破坏后产生的碎片所致。$^{[9][10][11]}$ 一些系外卫星可能具备成为地外生命栖息地的潜力。$^{[2]}$ 在双流浪行星系统 2MASS J11193254−1137466 AB 的其中一个组分周围，还发现了一颗可能的系外卫星；如果该天体真实存在，则可能具备宜居条件。$^{[12]}$
\subsection{定义与命名}
传统用法中，“卫星”通常指环绕行星运行的天体，但近年来发现的棕矮星及其行星尺度的伴随天体，使行星与卫星之间的界限变得模糊，其根源在于棕矮星本身质量较低。对此，国际天文学联合会（IAU）通过如下声明加以澄清：“凡真实质量低于氘热核聚变临界质量、并环绕恒星、棕矮星或恒星遗骸运行，且其与中心天体的质量比低于 L4/L5 不稳定性阈值（$M/M_{\text{central}} < 2/(25+\sqrt{621})$）的天体，定义为行星。”$^{[13]}$

IAU 的该定义并未涉及那些质量低于棕矮星、且低于氘燃烧极限的自由漂浮天体（这类天体通常被称为自由漂浮行星、流浪行星、低质量棕矮星或孤立行星质量天体）的卫星命名规范。在相关文献中，这些天体的卫星通常仍被称为系外卫星。$^{[7][8][12]}$

系外卫星的命名规则通常沿用其母体天体的名称，并在其后加上一个大写罗马数字。例如，Kepler-1625b 环绕 Kepler-1625（亦称 Kepler-1625a）运行，而 Kepler-1625b 本身可能拥有一颗名为 Kepler-1625b I 的卫星（目前尚未发现 Kepler-1625b II，也尚不清楚 Kepler-1625b I 是否拥有次级卫星）。
